\scanpage{104}
% archive foliation: 103
\deleted{\unsure{pro}}

\diagram{\(\begin{array}{c@{\qquad}ccc}
& \multicolumn{2}{c}{\overbrace{\hspace{9em}}^{\text{cas propre, ou du }H^1_!}} & \\
H^1 & \begin{array}{c}\text{poids }0\\\text{niveau }0\end{array}
& \begin{array}{c}\text{poids }1\\\text{niveau }1\end{array}
& \begin{array}{c}\text{poids }2\\\text{niveau }0\end{array}\\
& & \multicolumn{2}{c}{\underbrace{\hspace{9em}}_{\text{cas lisse}}}
\end{array}\)}

\margin{morceaux purs sont semi-simples}

\diagram{\(\begin{array}{c@{\qquad}c|ccccc}
& \multicolumn{1}{c}{} & \multicolumn{3}{c}{\overbrace{\hspace{13em}}^{\text{cas propre}}} & & \\
H^2_{\deleted{\unsure{!}}} & \text{poids} & 0 & 1 & 2 & 3 & 4\\
\cline{2-2}
& \text{niveau} & 0 & 1 & 2 & 1 & 0\\
& \multicolumn{1}{c}{} & & & \multicolumn{3}{c}{\underbrace{\hspace{13em}}_{\text{cas lisse}}}
\end{array}\)}

morceaux purs de poids $\ne2$ \unsure{st|sont} semi-simples.

\margin{$\Vert$}
Le morceau pur de poids $2$ est une extension par un motif de Tate $Q(-1)^r$.

$H^2_!$\qquad il semble qu'il puisse y avoir mélange de poids dans le mauvais ordre, à cause d'une extension d'un motif effectif de poids $0$ par un autre de poids $1$; en plus de cela, il \unsure{apparaît} une extension, peut-être non triviale, d'un motif effectif de poids $0$ \inserted{resp. $1$} par un autre.
